% QMB_n11_hardware.tex — Kapitel 11: Hardware
% Inhaltliche Grundlage: Dok. 083, 161, 183, 184; eigene Darstellung
\chapter{Hardware: Coherent Ising Machine, Google Willow, p-Bit, Photonenchip}
\label{ch:hardware}

\section{Coherent Ising Machine: Optimierung als physikalische Relaxation}

Die Coherent Ising Machine (CIM) ist ein optischer Analogcomputer, der
kombinatorische Optimierungsprobleme durch physikalische Relaxation in
den Grundzustand eines Ising-Hamiltonians löst. Das fundamentale Prinzip:
Das System \emph{fällt} in seine Lösung, es berechnet sie nicht.

\subsection*{Der Ising-Hamiltonian als universelle Optimierungssprache}

Jedes NP-harte kombinatorische Problem lässt sich als Minimierung des
Ising-Hamiltonians formulieren:
\begin{equation}
  H_{\mathrm{Ising}} = -\tfrac12\sum_{i\ne j}J_{ij}\sigma_i\sigma_j
  - \sum_i h_i\sigma_i,
\end{equation}
wobei $J_{ij}$ die Kopplungsstärken und $h_i$ externe Felder sind.
Der Grundzustand --- die Spin-Konfiguration mit minimalem $H_{\mathrm{Ising}}$
--- entspricht der optimalen Lösung. Die Abbildung ist exakt, keine
Approximation: Travelling Salesman, MaxCut, Portfolio-Optimierung,
Proteinstrukturvorhersage, RSA-Faktorisierung sind sämtlich als
Ising-Minimierung formulierbar.

\subsection*{Physikalische Implementierung}

Die CIM implementiert das Ising-Modell durch ein Netzwerk zeitgemulti-
plexter degenerierter optischer parametrischer Oszillatoren (DOPO):
Spin $+1$ ist ein Laserpuls mit Phase $0°$, Spin $-1$ ein Puls mit
Phase $180°$; die Kopplungen $J_{ij}$ werden durch elektronisches
Feedback realisiert. Das System entwickelt sich deterministisch gemäß
den Feldgleichungen des parametrischen Oszillators und findet dabei mit
hoher Wahrscheinlichkeit den Grundzustand. Bis zu 100\,000 optische
Spins sind realisiert (Stand 2024).\status{E}

Eine CIM scheint stochastisch zu suchen; verschiedene Läufe liefern
manchmal verschiedene Lösungen. Aber physikalisch folgt das System
\emph{deterministisch} den Feldgleichungen. Die scheinbare Zufälligkeit
ist ausschließlich Unkenntnis über die genauen Anfangsbedingungen des
Quantenvakuums --- eine direkte Parallele zur FFGFT-Messung.

\subsection*{Sechs Strukturelle Parallelen zur FFGFT}

Zwischen CIM und FFGFT lassen sich sechs Parallelen identifizieren:

\textbf{Parallele 1: Determinismus hinter scheinbarer Stochastik.} In
der CIM ist die Lösung physikalisch erzwungen; in der FFGFT ist das
Messergebnis durch die Feldkonfiguration geometrisch bestimmt. Beide
ersetzen fundamentalen Probabilismus durch versteckten Determinismus.

\textbf{Parallele 2: Grundzustand als geometrisches Minimum.} Die CIM
relaxiert in das Energieminimum des Ising-Hamiltonians; die FFGFT
beschreibt das Universum als im Grundzustand seiner eigenen Geometrie
mit $\xi = 4/30000$. Optimierung ist Relaxation, nicht Iteration.

\textbf{Parallele 3: Qubit-Struktur.} CIM-Spin $\pm1$ ist die Phase
des Laserpulses; FFGFT-Qubit $|\tilde T_\uparrow,m_\downarrow\rangle$
und $|\tilde T_\downarrow,m_\uparrow\rangle$ ist durch $\tilde T\cdot m = 1$
intrinsisch erzwungen. Der FFGFT-Spin muss nicht extern auf $\pm1$
begrenzt werden; er ist es durch die Geometrie.

\textbf{Parallele 4: Die Kopplungsmatrix.} In der klassischen CIM sind
$J_{ij}$ extern programmiert; in einer FFGFT-basierten Ising-Maschine
würden sie geometrisch aus dem Zeitfeld emergieren:
$J_{ij}^{(\mathrm{FFGFT})} = \int\tilde T_i\cdot\tilde T_j\cdot G(x_i,x_j)\,d^3x$.
Das bedeutet: keine Kodierungsfehler möglich, da die Kopplungen reale
physikalische Wechselwirkungen wären.\status{S}

\textbf{Parallele 5: Fraktale Korrektur als Annealing.} Der Korrekturfaktor
$K_{\mathrm{frak}}^{3/2}$ spielt in der FFGFT die analoge Rolle des
Annealing-Temperaturparameters: Er verhindert, dass das System in einem
lokalen Minimum der geometrischen Energie verbleibt.

\textbf{Parallele 6: Toroidale Topologie.} Die FFGFT beschreibt den
Träger als $T^4/\mathbb{Z}_3$. Toroidale Randbedingungen hätten für
eine physikalische Ising-Maschine einen Vorteil: kein Randproblem, alle
Spins haben gleich viele Nachbarn, natürliche Implementierung
translationsinvarianter Kopplungen.

\section{Google Willow: topologische Fehlerkorrektur und Selbstmessung}

Googles Willow-Quantenprozessor erzielte 2024 und 2025 drei Durchbrüche,
die aus FFGFT-Sicht besondere Bedeutung haben.

\subsection*{Fehlerkorrektur unter dem Schwellenwert (Dezember 2024)}

Willow (105 supraleitende Transmon-Qubits, $\sim10$\,mK) zeigt erstmals,
dass die logische Fehlerrate exponentiell sinkt, wenn die Code-Distanz
$d$ erhöht wird.\status{E} Ein $d = 7$-Block erzielt eine Fehlerrate um
Faktor $2{,}14$ unter dem $d = 5$-Block. Physikalische Einzel-Qubit-
Fehlerraten liegen bei $0{,}05\,\%$ --- halb so viel wie auf dem
Vorgänger Sycamore. Die Fehlerkorrektur basiert auf einem Surface Code
mit topologisch geschützten logischen Qubits.

FFGFT-Einordnung: Mehr physische Qubits reduzieren den logischen Fehler
exponentiell --- das ist der experimentelle Beleg, dass topologisch stabile
Konfigurationen im Quantensystem existieren und erreichbar sind, nicht
durch externe Rauschunterdrückung, sondern durch geometrische Ordnung.
Das ist die Richtung, auf die die FFGFT mit ihren Modenklassen des
$D_4$-Gitters zeigt (Kapitel~\ref{ch:gitter}).\status{S}

\subsection*{Floquet topologische Ordnung (September 2025)}

Will et al.\ realisierten auf Willow eine Floquet-topologisch geordnete
Phase --- eine Quantenphase, die im thermischen Gleichgewicht verboten
ist, aber unter periodischem äußerem Antrieb entsteht.\status{E} Auf bis
zu 58 Qubits entstanden: chirale Kantenmoden (Informationsausbreitung in
ausgezeichneter Richtung), anyonische Anregungen (e-Anyonen werden zu
m-Anyonen und umgekehrt), und ein interferometrisch messbarer
topologischer Bulk-Invariant.

FFGFT-Einordnung: Die nicht-gleichgewichtliche Ordnung trägt eine
$\mathbb{Z}_2$-Symmetrie auf dem zweidimensionalen Gitter --- strukturell
verwandt mit der $\mathbb{Z}_3$-Trialität des FFGFT-Trägers, aber nicht
identisch (Dimension, Symmetriegruppe und Trägergeometrie verschieden).
Die Floquet-topologische Ordnung ist ein experimentelles Analogon der
Modenstruktur auf $T^4/\mathbb{Z}_3$.\status{S}

\subsection*{Quantum Echoes (Oktober 2025)}

Out-of-Time-Order-Korrelatoren (OTOC) auf Willow zeigen harmonische
Ausbreitung von Quanteninformation durch konstruktive Interferenz --- ein
Echo-Muster, das periodisch zurückkehrt. FFGFT-Einordnung: Im
Torusformalismus sind Moden auf einer geschlossenen Geometrie
definiert; Informationsausbreitung kehrt zur Quelle zurück --- das Echo
ist eine Konsequenz der Kompaktheit, nicht ein Zufall.

\section{Das p-Bit: zwischen Bit und Qubit}

Das p-Bit (probabilistisches Bit) besetzt den Zwischenbereich zwischen
dem klassischen Bit (tiefe Potenzialsenke, $U\gg k_BT$, stabil über
Jahre) und dem Qubit (metastabile Superposition, $U\ll k_BT$, zerfällt in
Pikosekunden bis Millisekunden). Es ist eine magnetische Nanostruktur auf
der geometrischen Skala $L_8\approx22$\,nm, bei der die Torsionsbarriere
thermisch überwindbar wird.

\subsection*{Physikalische Grundlage}

Die Energiebarriere ist die Stoner-Wohlfarth-Energie $U = K_u\cdot V$,
mit der magnetischen Anisotropieenergie $K_u$ und dem Volumen $V$. Das
Schaltverhalten folgt dem Néel-Arrhenius-Gesetz:
$\tau = \tau_0\exp(U/k_BT)$. Bei 22\,nm und Raumtemperatur liegt $\tau$
im Bereich von Mikrosekunden bis Millisekunden --- kein stabiles Bit,
aber auch kein kohärentes Qubit.

Die FFGFT-Einordnung: Das Schaltverhalten ist deterministisch auf der
sub-Planck-Zeitskala $t_0 = t_P/7500\approx3{,}6\times10^{-47}$\,s;
was als thermische Stochastizität erscheint, ist die makroskopisch
unaufgelöste Abtastung dieser deterministischen Dynamik. Die scheinbare
Zufälligkeit ist epistemisch, nicht ontologisch.\status{S}

\subsection*{Rekonfigurierbare Funktion}

Durch Wahl der magnetischen Vorspannung kann dasselbe p-Bit als
stochastischer Schalter, als erzwungener Schalter oder als analoger
Komparator betrieben werden. Das macht p-Bits zu Bausteinen für
probabilistische Schaltkreise (p-Bits-Netzwerke), die wie Boltzmann-
Maschinen trainiert werden können und NP-harte Probleme ähnlich wie CIMs
durch physikalische Relaxation angehen.

FFGFT-Einordnung: Das p-Bit ist kein Quantenobjekt (kein kohärenter
Überlagerungszustand), sondern eine thermisch fluktuierende klassische
Mode. Im Torusformalismus entspricht es einer Mode, deren Phase $\theta$
durch thermisches Rauschen vollständig randomisiert ist --- der Grenzfall
maximaler Dekohärenz auf dem Träger.

\subsection*{Warum 22 nm: die geometrische Schwelle}

Die kritische Größe für ein p-Bit bei Raumtemperatur entspricht der
FFGFT-Stufe $N = 8{,}00$, also $L_8\approx22$\,nm. Das ist dieselbe
Schwelle, bei der klassische Transistoren 2012 die Quantengrenze
erreichten (Kapitel~\ref{ch:spin}). Drei unabhängige Phänomene konvergieren
auf dieselbe geometrische Skala: Transistorgrenze, exzitonischer
Bohr-Radius, p-Bit-Schwelle.\status{S}

\section{Photonenchip: Raumtemperatur-Quantenphotonik}

Chinas CHIPX/Touring-Quantum-Chip (2025) ist ein monolithischer
6-Zoll-Dünnfilm-Lithiumniobat (TFLN)-Wafer mit über 1000 optischen
Komponenten. Kerndaten: $1000\times$-Speedup für parallele Tasks,
$100\times$ geringerer Energieverbrauch, stabil bei Raumtemperatur,
$12000$ Wafer pro Jahr Produktion.\status{E}

\textbf{Warum Raumtemperatur funktioniert:} Optische Photonen haben
$E\approx1$--2\,eV; $E/k_BT(300\,\text{K})\approx40$. Thermisches
Rauschen kann keine optischen Photonen erzeugen; das System sitzt weit
unterhalb der thermischen Schwelle (Kapitel~\ref{ch:qubit}). Der TFLN-
Einkristall mit $\sim10^{23}$ Gitterplätzen ist geometrisch extrem
robust.\status{K}

\textbf{FFGFT-Einordnung:} Masselose Moden (Photonen) sind in der FFGFT
Dreier-Orbits in $GF(27)^*$ --- strukturell von den massiven Fermionen
(Fixpunkte $\{+1,-1\}$) getrennt (Frobenius-Trennung, Dok.~339).\status{B}
Die photonische Plattform eignet sich für schlupflochfreie Bell-Tests und
ist damit der natürliche Prüfstand für die $\xi$-Abweichung
$\Delta\mathrm{CHSH}\approx\xi/(2\pi)\approx2\times10^{-5}$ --- wenn
Detektionseffizienz und Kohärenzzeiten ausreichen.

\textbf{Vorgeschlagene FFGFT-Optimierungen:} Drei Ansätze sind konzeptuell
skizziert, alle noch nicht realisiert: (1) Topogie-Compiler, der
topologische Qubit-Platzierung nach der $D_4$-Gitterstruktur vornimmt
und SWAP-Gatter um 30--50\,\% reduzieren könnte; (2) harmonische Resonanz
mit Qubit-Frequenzen auf $\xi$-Vielfachen; (3) Zeitfeld-Modulation zur
aktiven Kohärenzerhaltung. Status: [S] für alle drei.

\quelle{Dok.~083 (Photonenchip 2025), Dok.~161 (Coherent Ising Machine), Dok.~183 (Google Willow, April 2026), Dok.~184 (p-Bit), Dok.~339 (Frobenius-Trennung), Dok.~343 (Auflösungsboden)}
